\section{均匀电子液（jellium）哈密顿量}

本节按 Giuliani--Vignale（以下简称 GV）《Quantum Theory of the Electron Liquid》第 1 章，给出 $d=2,3$ 维均匀电子液的设定、标度与二次量子化哈密顿量。

\subsection{模型设定与库仑正则化}

体积 $L^d$ 内 $N$ 个电子，正电荷背景密度 $+en$ 以保证电中性（jellium）。$d$ 维库仑作用的 Fourier 分量为
\begin{equation}
  v_{\mathbf{q}}
  =
  \begin{cases}
    4\pi e^2/q^2, & d=3,\\
    2\pi e^2/q, & d=2.
  \end{cases}
\label{eq:vq_d}
\end{equation}
$q\to0$ 发散由电子--电子、电子--背景、背景--背景三项相消。正则化后
\begin{equation}
  H=\sum_{\mathbf{k}\sigma}
  \frac{\hbar^2k^2}{2m}\,
  c_{\mathbf{k}\sigma}^\dagger c_{\mathbf{k}\sigma}
  +\frac{1}{2L^d}\sum_{\mathbf{q}\neq0}
  v_{\mathbf{q}}
  \bigl(\hat n_{\mathbf{q}}\hat n_{-\mathbf{q}}-\hat N\bigr).
\label{eq:jelliumH}
\end{equation}
一维需额外短距截断，因 $v_q\propto -\ln(qa)$ 不再齐次，标度分析与 $d=2,3$ 不同。

\subsection{密度参数 \texorpdfstring{$r_s$}{rs} 与能量单位}

Wigner--Seitz 半径由每个电子占据的 $d$ 维球体积定义：
\begin{equation}
  n=\frac{N}{L^d}
  =
  \begin{cases}
    \bigl[\tfrac43\pi (r_s a_B)^3\bigr]^{-1}, & d=3,\\
    \bigl[\pi (r_s a_B)^2\bigr]^{-1}, & d=2,
  \end{cases}
  \qquad
  a_B=\frac{\hbar^2}{me^2}.
\end{equation}
费米波矢
\begin{equation}
  k_F
  =
  \begin{cases}
    (3\pi^2 n)^{1/3}, & d=3\text{（顺磁）},\\
    (2\pi n)^{1/2}, & d=2\text{（顺磁）}.
  \end{cases}
\end{equation}
以 Bohr 半径为长度、Rydberg $\mathrm{Ry}=e^2/(2a_B)$ 为能量，无量纲化后 $H$ 只依赖 $r_s$：动能 $\propto r_s^{-2}$，库仑 $\propto r_s^{-1}$。故 $r_s\ll1$ 为高密弱耦合，$r_s\gg1$ 为低密强耦合。

半导体二维气中用有效质量 $m_b$ 与背景介电常数 $\varepsilon_b$ 定义
\begin{equation}
  a_B^*=\frac{\hbar^2\varepsilon_b}{m_b e^2},
  \qquad
  \mathrm{Ry}^*=\frac{e^2}{2\varepsilon_b a_B^*}.
\end{equation}
数量级换算（GV 表 1.2）：三维 $n\sim10^{22}\,\mathrm{cm}^{-3}$ 对应 $r_s\sim5.45\,m_b/(\bar n^{1/3}\varepsilon_b)$；二维 $n\sim10^{11}\,\mathrm{cm}^{-2}$ 对应 $r_s\sim337\,m_b/(\bar n^{1/2}\varepsilon_b)$。

\subsection{无相互作用气与交换能（一阶）}

顺磁无相互作用每粒子动能
\begin{equation}
  \varepsilon_0(r_s)
  =\frac{A_d}{r_s^2}\,\mathrm{Ry},
  \qquad
  A_3=\frac35\Bigl(\frac{9\pi}{4}\Bigr)^{2/3}\approx2.21,
  \qquad
  A_2=1.
\end{equation}
平面波 Slater 行列式的一阶库仑能即交换能
\begin{equation}
  \varepsilon_x(r_s)
  =-\frac{B_d}{r_s}\,\mathrm{Ry},
  \qquad
  B_3=\frac{3}{2\pi}\Bigl(\frac{9\pi}{4}\Bigr)^{1/3}\approx0.916,
  \qquad
  B_2=\frac{8\sqrt{2}}{3\pi}\approx1.20.
\label{eq:ex_Bd}
\end{equation}
从而
\begin{equation}
  \varepsilon
  =\frac{A_d}{r_s^2}-\frac{B_d}{r_s}+\varepsilon_{\mathrm{c}}(r_s)
  \quad(\mathrm{Ry}).
\label{eq:eps_rs}
\end{equation}
$\varepsilon_{\mathrm{c}}$ 为关联能，需 RPA、耦合常数积分或 QMC。高密极限 $\varepsilon_{\mathrm{c}}$ 在三维以 $a\ln r_s+b+\cdots$ 展开（Gell-Mann--Brueckner）。
